\chapter{Line, Surface, and Domain-Wall Operators}
\label{ch:global-line-surface-domain-wall-operators}

\ConventionsMixed

Extended operators are QFT data supported on submanifolds.  They are not
determined by local operators alone.  This chapter defines the basic objects
needed for later discussions of confinement, screening, duality, and anomaly
inflow.

\section{Defect Data}

Let \(X\) be a spacetime manifold and let \(Y\subset X\) be an embedded
submanifold.  A defect supported on \(Y\) is specified by:
\begin{enumerate}[label=(\roman*)]
\item a choice of fields or boundary conditions on the complement
  \(X\setminus Y\);
\item singular behavior or extra degrees of freedom along \(Y\);
\item a coupling between bulk and defect variables;
\item a renormalization prescription for correlation functions with the
  defect inserted;
\item an equivalence relation identifying defect representatives with the
  same correlation functionals.
\end{enumerate}
The codimension of \(Y\) determines whether the object is a local operator,
line, surface, domain wall, boundary condition, or higher defect.

\section{Wilson Lines}

Let \(G\) be a compact gauge group and \(R\) a finite-dimensional
representation.  For an oriented curve \(\gamma\), the Wilson line is the
parallel-transport functional
\[
  W_R(\gamma)=
  \operatorname{Tr}_R\,\operatorname{Pexp}
  \int_\gamma A .
\]
In a regulated quantum gauge theory this expression is a coordinate for a
renormalized line operator.  Its definition requires the global form of
\(G\), the representation \(R\), framing or endpoint data when needed, and a
renormalization scheme for the short-distance singularity along the line.
If \(\gamma\) is open, the endpoints must carry charged states or additional
operators so that the full insertion is gauge invariant.

\section{Magnetic and Dyonic Lines}

A magnetic line in four-dimensional gauge theory is defined by a prescribed
singularity of the gauge field near a curve.  Around a small sphere linking
the line, the magnetic flux is labelled by a cocharacter
\[
  m:U(1)\longrightarrow G
\]
modulo conjugation, or equivalently by an element of the magnetic cocharacter
lattice modulo the Weyl group.  A dyonic line carries both the magnetic label
\(m\) and an electric representation of the stabilizer of \(m\).  The allowed
electric and magnetic labels obey the Dirac pairing condition: the phase
obtained by moving one line around another must be trivial for mutually local
operators in the chosen theory.

The global form of the gauge group changes the allowed labels.  For example,
the electric lattice of \(SU(N)\) differs from that of \(PSU(N)\), and the
dual magnetic lattice changes accordingly.  Therefore the line-operator
spectrum is part of the definition of a gauge theory: it is fixed by the
global gauge group, the allowed bundles, and the chosen mutually local
electric-magnetic charge lattice.

\begin{definition}[Renormalized 't Hooft line]
\label{def:renormalized-thooft-line}
Let \(X\) be an oriented four-manifold, let \(\gamma\subset X\) be an oriented
embedded curve, and let \(G\) be a compact connected gauge group with maximal
torus \(T\) and Weyl group \(W\).  A bare magnetic line of magnetic charge
\[
  m\in\operatorname{Hom}(U(1),T)/W
\]
is defined by removing a tubular neighborhood \(N_\epsilon(\gamma)\) and
integrating over \(G\)-connections on \(X\setminus N_\epsilon(\gamma)\) whose
restriction to each small linking sphere \(S^2_p\) lies in the topological
class determined by \(m\).  In a local trivialization over
\(\gamma\times(S^2_p\setminus\{\text{Dirac ray}\})\), this means that the
connection is gauge-equivalent near the boundary to the embedded Dirac
monopole
\[
  F\sim \frac{m_*}{2}\sin\theta\,\dd\theta\wedge\dd\varphi ,
  \qquad
  \int_{S^2_p}\frac{F}{2\pi}=m_* ,
\]
where \(m_*\in\mathfrak t\) is the differential of the cocharacter in the
chosen normalization.  Gauge transformations are required to preserve the
cocharacter class on the boundary.  Let \(\mathcal A_m\) denote the resulting
space of regulated connections and let \(\mathcal G_m\) denote the
corresponding boundary-preserving gauge group.  The renormalized 't Hooft line
\(H_m(\gamma)\) is the limit
\[
  H_m(\gamma)
  =
  \lim_{\epsilon\to0}
  Z_{\rm line}(\epsilon)\,
  \int_{\mathcal A_m(X\setminus N_\epsilon(\gamma))/\mathcal G_m}
  \mathcal D A\,\exp(-S[A])
\]
as a correlation functional with the same additional insertions, after
multiplication by local line counterterms \(Z_{\rm line}(\epsilon)\).  The
definition exists only in a specified regulator or constructive framework in
which the limit and the allowed counterterms are part of the line-operator
data.
\end{definition}

For \(G=U(1)\), the cocharacter lattice is \(\mathbb Z\).  The minimal
't Hooft line \(H(\gamma)\) has magnetic charge \(1\), so the boundary
condition is
\[
  \int_{S^2_p}\frac{F}{2\pi}=1
\]
for every small linking sphere.  If the line also carries electric charge
\(q\), one multiplies the magnetic boundary condition by a Wilson factor for
the unbroken gauge group along \(\gamma\); this gives a dyonic line.  A
dynamical monopole operator is an endpoint for such a magnetic line.  Hence a
theory containing dynamical monopoles does not have the magnetic one-form
symmetry under which closed 't Hooft lines are charged.

\section{Surface Operators}

A surface operator supported on a two-dimensional submanifold \(\Sigma\) may
be defined by singular behavior of the gauge field, by coupling a two-
dimensional QFT on \(\Sigma\) to the bulk, or by inserting a higher-form
charged object.  In four-dimensional gauge theory an abelianized surface
singularity has local model
\[
  A\sim \alpha\, d\varphi
\]
near the normal plane to \(\Sigma\), where \(\alpha\) is valued in a Cartan
subalgebra modulo the cocharacter lattice and Weyl group.  Additional
parameters can couple to the magnetic flux through \(\Sigma\).  A complete
definition must state the singularity, defect counterterms, gauge
equivalence, and how local operators approaching the surface are
renormalized.

\section{Domain Walls and Interfaces}

A domain wall between theories \(T_-\) and \(T_+\) is a codimension-one defect
whose correlation functional takes bulk fields of \(T_-\) on one side and
bulk fields of \(T_+\) on the other.  A topological wall is one whose
correlators are invariant under deformations of the wall that avoid other
insertions.  Such a wall acts on local and extended operators by bringing an
operator to the wall and moving it through, provided the required limits
exist.

\begin{definition}[Fusion of topological defects]
\label{def:fusion-topological-defects}
Let \(D_1\) and \(D_2\) be parallel topological defects of the same
codimension.  Their fusion \(D_1\otimes D_2\), when it exists, is the defect
defined by the limit in which the separation between \(D_1\) and \(D_2\) goes
to zero, after subtracting or renormalizing local defect counterterms.
\end{definition}

Fusion is an operation only after the limiting procedure and equivalence
relation have been supplied.  In favorable topological sectors the resulting
objects form a category or higher category; in a general QFT the short-distance
limit can produce new defect degrees of freedom or fail to exist.

\begin{openproblem}[Extended-operator reconstruction]
\label{op:extended-operator-reconstruction}
Given the local operator algebra and symmetry data of a QFT, determine what
additional input is required to reconstruct its extended operators.  The
answer must include global forms, higher-form symmetries, singular boundary
conditions, defect renormalization, and the topology in which defect fusion
limits are taken.
\end{openproblem}
